% Pursuing Stacks (cotes 134-2 to 134-8): archive cote 134-6, site page 121 (PDF page 122), typescript page 439.
% typescript. Draft compared: ps5.tex:7854-7918.
% AI-transcribed by Claude Opus 5.5 (claude-opus-5-5, adaptive thinking, effort high), Message Batches API, batch ps3 (msgbatch_01XdFmuVVBAcgc7EFxdtCqsj), 2026-11-12; not yet checked by eye.
\scanpage{134-6/121}
\mspage{439}

Q = P$^{\mathrm{O}}$). We'll come back upon this later, I expect.  This structure
($\ast$) will be the k-linear analogon of the usual diagonal map
\[ (\ast\ast) \qquad A \hand{$\to$} A\mathrm{x}A \]
for a small category A,  giving rise by k-linearization to
\[ \mathrm{Add}_k(A) \hand{$\to$} \mathrm{Add}_k(A\mathrm{x}A) \hand{$\unsure{\approxeq}$} \mathrm{Add}_k(A)\otimes_k\mathrm{Add}_k(A) \quad , \]
namely a structure of type ($\ast$).  It just occured to me,  through the
reflections of these last days,  that the ``coalgebra structure'' ($\ast$) may
well turn out \inserted{(taking $k = \unsure{Z}$)} to be the more sophisticated structure than a usual coal-%
gebra structure (cf. p.339  (27)), \inserted{needed} in order to grasp ``in linear terms''
the notion of a homotopy type,  possibly under restrictions such as
1-connectedness, .as pondered about in section 94.  This looks at any
rate a more ``natural'' object than the De Rham complex with devided
powers\xout{t}, referred to in loc. cit. ,  and is more evidently adapted to
our point of view of using small categories as models for homotopy
types.  The greater sophistication,  in comparison to De Rham type com-%
plexes, lies in this,  that here the objects serving as models (whether
small categories,  or small additive categories endowed with \inserted{a} diagonal
map) are objects in a \uline{2-category},  \over{\unsure{w}}{w}hereas De Rham complexes and the \xout{lik}
like are just objects in ordinary categories,  without any question of
taking ``maps between maps''.  This feature implies, ``as usual'' (or in
duality rather to familiar situations with tensor product functors ...),
that the \inserted{(anti)} commutative and associative \uline{axioms}  familiar from linear alge-%
bra\hand{,} \hand{(}in the case of usual k-algebras or k-coalgebras\hand{)}\hand{,} should be
replaced by commutativity and associativity \uline{data},  namely given isomor-%
phisms (not identities) between \xout{the} two natural functors
\[ P \hand{$\to$} P\otimes_k P \quad , \quad P \hand{$\longrightarrow$} P\otimes_k P\otimes_k P \]
deduced from ($\ast$).  The axioms now will be more sophisticated,  they will
express ``coherence conditions'' on these data  -  one place maybe where
this is developped somewhat,  in the context of diagonal maps ($\ast$),
might be Saavedra's thesis.  (The  more familiar case,  when starting
with a tensor product operation on a category,  together with associati-%
vity and \inserted{/or} commutativity data,  has been done with \xout{\unsure{some}} care by various
mathematicians,  including Mac Lane, Benabou,  Mme Sinh Hoang Xuan,
and presumably it should be enough to ``reverse arrows'' in order to get
``the'' natural set of coherence axioms for a diagonal map ($\ast$)\hand{)}.  The
``intriguing feature''  with the would-be De Rham models for homotopy
types (cf. p. 341, 342) ,  namely that the latter make a sense over any

\margin{121}
